% Options for packages loaded elsewhere
\PassOptionsToPackage{unicode}{hyperref}
\PassOptionsToPackage{hyphens}{url}
\documentclass[
  a4paper,
]{article}
\usepackage{xcolor}
\usepackage[margin=0.75in]{geometry}
\usepackage{amsmath,amssymb}
\setcounter{secnumdepth}{3} % number main sections and subsections
\usepackage{iftex}
\ifPDFTeX
  \usepackage[T1]{fontenc}
  \usepackage[utf8]{inputenc}
  \usepackage{textcomp} % provide euro and other symbols
\else % if luatex or xetex
  \usepackage{unicode-math} % this also loads fontspec
  \defaultfontfeatures{Scale=MatchLowercase}
  \defaultfontfeatures[\rmfamily]{Ligatures=TeX,Scale=1}
\fi
\usepackage{lmodern}
\ifPDFTeX\else
  % xetex/luatex font selection
\fi
% Use upquote if available, for straight quotes in verbatim environments
\IfFileExists{upquote.sty}{\usepackage{upquote}}{}
\IfFileExists{microtype.sty}{% use microtype if available
  \usepackage[]{microtype}
  \UseMicrotypeSet[protrusion]{basicmath} % disable protrusion for tt fonts
}{}
\makeatletter
\@ifundefined{KOMAClassName}{% if non-KOMA class
  \IfFileExists{parskip.sty}{%
    \usepackage{parskip}
  }{% else
    \setlength{\parindent}{0pt}
    \setlength{\parskip}{6pt plus 2pt minus 1pt}}
}{% if KOMA class
  \KOMAoptions{parskip=half}}
\makeatother
\usepackage{color}
\usepackage{fancyvrb}
\newcommand{\VerbBar}{|}
\newcommand{\VERB}{\Verb[commandchars=\\\{\}]}
\DefineVerbatimEnvironment{Highlighting}{Verbatim}{commandchars=\\\{\}}
% Add ',fontsize=\small' for more characters per line
\newenvironment{Shaded}{}{}
\newcommand{\AlertTok}[1]{\textcolor[rgb]{1.00,0.00,0.00}{\textbf{#1}}}
\newcommand{\AnnotationTok}[1]{\textcolor[rgb]{0.38,0.63,0.69}{\textbf{\textit{#1}}}}
\newcommand{\AttributeTok}[1]{\textcolor[rgb]{0.49,0.56,0.16}{#1}}
\newcommand{\BaseNTok}[1]{\textcolor[rgb]{0.25,0.63,0.44}{#1}}
\newcommand{\BuiltInTok}[1]{\textcolor[rgb]{0.00,0.50,0.00}{#1}}
\newcommand{\CharTok}[1]{\textcolor[rgb]{0.25,0.44,0.63}{#1}}
\newcommand{\CommentTok}[1]{\textcolor[rgb]{0.38,0.63,0.69}{\textit{#1}}}
\newcommand{\CommentVarTok}[1]{\textcolor[rgb]{0.38,0.63,0.69}{\textbf{\textit{#1}}}}
\newcommand{\ConstantTok}[1]{\textcolor[rgb]{0.53,0.00,0.00}{#1}}
\newcommand{\ControlFlowTok}[1]{\textcolor[rgb]{0.00,0.44,0.13}{\textbf{#1}}}
\newcommand{\DataTypeTok}[1]{\textcolor[rgb]{0.56,0.13,0.00}{#1}}
\newcommand{\DecValTok}[1]{\textcolor[rgb]{0.25,0.63,0.44}{#1}}
\newcommand{\DocumentationTok}[1]{\textcolor[rgb]{0.73,0.13,0.13}{\textit{#1}}}
\newcommand{\ErrorTok}[1]{\textcolor[rgb]{1.00,0.00,0.00}{\textbf{#1}}}
\newcommand{\ExtensionTok}[1]{#1}
\newcommand{\FloatTok}[1]{\textcolor[rgb]{0.25,0.63,0.44}{#1}}
\newcommand{\FunctionTok}[1]{\textcolor[rgb]{0.02,0.16,0.49}{#1}}
\newcommand{\ImportTok}[1]{\textcolor[rgb]{0.00,0.50,0.00}{\textbf{#1}}}
\newcommand{\InformationTok}[1]{\textcolor[rgb]{0.38,0.63,0.69}{\textbf{\textit{#1}}}}
\newcommand{\KeywordTok}[1]{\textcolor[rgb]{0.00,0.44,0.13}{\textbf{#1}}}
\newcommand{\NormalTok}[1]{#1}
\newcommand{\OperatorTok}[1]{\textcolor[rgb]{0.40,0.40,0.40}{#1}}
\newcommand{\OtherTok}[1]{\textcolor[rgb]{0.00,0.44,0.13}{#1}}
\newcommand{\PreprocessorTok}[1]{\textcolor[rgb]{0.74,0.48,0.00}{#1}}
\newcommand{\RegionMarkerTok}[1]{#1}
\newcommand{\SpecialCharTok}[1]{\textcolor[rgb]{0.25,0.44,0.63}{#1}}
\newcommand{\SpecialStringTok}[1]{\textcolor[rgb]{0.73,0.40,0.53}{#1}}
\newcommand{\StringTok}[1]{\textcolor[rgb]{0.25,0.44,0.63}{#1}}
\newcommand{\VariableTok}[1]{\textcolor[rgb]{0.10,0.09,0.49}{#1}}
\newcommand{\VerbatimStringTok}[1]{\textcolor[rgb]{0.25,0.44,0.63}{#1}}
\newcommand{\WarningTok}[1]{\textcolor[rgb]{0.38,0.63,0.69}{\textbf{\textit{#1}}}}
\usepackage{longtable,booktabs,array}
\usepackage{caption}
\newcounter{none} % for unnumbered tables
\usepackage{calc} % for calculating minipage widths
% Correct order of tables after \paragraph or \subparagraph
\usepackage{etoolbox}
\makeatletter
\patchcmd\longtable{\par}{\if@noskipsec\mbox{}\fi\par}{}{}
\makeatother
\setlength{\emergencystretch}{3em} % prevent overfull lines
\providecommand{\tightlist}{%
  \setlength{\itemsep}{0pt}\setlength{\parskip}{0pt}}
% \DeclareUnicodeCharacter is available only with the pdfTeX inputenc path.
% Unicode engines handle these characters natively via fontspec/unicode-math.
\ifPDFTeX
\DeclareUnicodeCharacter{2264}{\ensuremath{\leq}}
\DeclareUnicodeCharacter{2265}{\ensuremath{\geq}}
\DeclareUnicodeCharacter{2248}{\ensuremath{\approx}}
\DeclareUnicodeCharacter{2212}{\ensuremath{-}}
\DeclareUnicodeCharacter{26A0}{\textbf{[!]}}
\DeclareUnicodeCharacter{FE0F}{}
\fi
\usepackage{bookmark}
\IfFileExists{xurl.sty}{\usepackage{xurl}}{} % add URL line breaks if available
\urlstyle{same}
% fallback for those not using the hyperref driver hyperxmp:
\makeatletter
\@ifundefined{xmpquote}{\newcommand{\xmpquote}[1]{#1}}{}
\makeatother
\hypersetup{
  hidelinks,
  pdfcreator={LaTeX}}

\title{Agent Sandbox Lifecycle Churn: A Use-case-grounded Benchmark Proposal}
\author{Bingrui Jin \qquad Natalia Naranjo Rodríguez\\[0.5em]\small Supervisor: Chih-Kai Huang}
\date{Télécom Paris --- November 2026}

\begin{document}
\maketitle
\setcounter{tocdepth}{2}
\tableofcontents

\section{Introduction and Research Gap}
\label{sec:introduction}
\subsection{Motivation and research question}
Sandbox startup is only one event in an agent execution
environment's lifecycle. Coding agents issue multiple
commands, pause for review and reuse state; evaluation
coordinators create and release many independent trials;
and persistent agents may receive new requests after
long inactivity. These application events generate
provisioning, scheduling, execution, retention and
cleanup work on Kubernetes and platform-specific
controllers or worker pools.

The central research question is:
\begin{quote}
\textbf{How do Kubernetes-based sandbox architectures
handle lifecycle churn, and where do control-plane
and runtime bottlenecks emerge under sustained,
bursty and intermittent agent workloads?}
\end{quote}

Rather than rank sandbox products by a single
ready timestamp, we follow a user-level request
through the first independently verified command
or task result, together with the infrastructure
operations required to deliver it.

\subsection{Existing work and the research gap}
Prior research and engineering benchmarks address
important parts of sandbox performance, but typically
under different workloads, hardware, isolation runtimes
or state contracts. Their results should not be
combined into an uncontrolled leaderboard.

{\small
\begin{longtable}{@{}p{0.24\linewidth}p{0.31\linewidth}p{0.37\linewidth}@{}}
\toprule
Prior work & Established focus & Comparative question left open\\
\midrule
\endhead
\bottomrule
\endlastfoot
Sandman \cite{sandman} &
High-concurrency provisioning and lifecycle
management for production agentic RL. &
How do different Kubernetes-managed architectures
behave for the same independently verified task?\\
SIGs Agent Sandbox tuning \cite{sigTuning} &
Controller concurrency, warm claims,
bursts, sustained demand and API overhead. &
Where do lifecycle costs move between
Kubernetes and specialized controllers/workers?\\
ARIES \cite{ariesPaper,ariesRepo} &
Agent trajectories, runtime experiments
and resource attribution. &
How do real task events map to native
sandbox lifecycle work and state guarantees?\\
Checkpoint and Fast Sandbox studies
\cite{deltaBox,fastPerformance} &
Snapshot/restore performance, prepared
templates and runtime-specific paths. &
How do preparation, state fidelity and
client-visible execution affect comparisons?\\
Open benchmarks
\cite{frigaBenchmark,computeBenchmark} &
Hosted tool/first-command timings and
independent execution checks. &
What happens under matched Kubernetes
hardware with observability across layers?\\
\end{longtable}
}

The intended gap is a \emph{cross-architecture,
use-case-grounded lifecycle benchmark} with a
common task contract, correctness checks and
qualified attribution across Kubernetes,
specialized management components and execution
runtimes. This is a proposed research direction,
not an unverified claim of absolute novelty.

\subsection{Research questions and intended contributions}
Three questions guide the experiment.
\begin{enumerate}
\item \textbf{RQ1 (cost placement):} Which Kubernetes,
specialized control-plane, runtime and storage
operations does each workload induce?
\item \textbf{RQ2 (load response):} Where do
verified throughput, latency and reliability
degrade under bursts and sustained churn?
\item \textbf{RQ3 (idle/state):} How do native
idle-management policies change user-return
latency, resource cost and state correctness?
\end{enumerate}

Subject to successful deployment and measurements,
potential contributions are a reproducible
application-derived benchmark, a state-aware
comparison, cross-layer bottleneck evidence,
and conditional architecture--workload trade-offs.
The study is not designed to declare a universal
best sandbox platform.

\section{Systems and Architecture Taxonomy}
\label{sec:systems}
\subsection{Kubernetes-based candidate architectures}
Five candidate configurations remain in scope.
They do not all map one logical sandbox to one
Pod: some use Kubernetes-managed pools of
workers that admit further guest or actor execution.

{\small
\begin{longtable}{@{}p{0.24\linewidth}p{0.45\linewidth}p{0.23\linewidth}@{}}
\toprule
Candidate & Lifecycle architecture and distinction & Feasibility priority\\
\midrule
\endhead
\bottomrule
\endlastfoot
Native Kubernetes Pod &
Direct Pod scheduling, execution and deletion.
Cold baseline; an optional pre-created warm-Pod
allocator is a \emph{separate configuration}. &
P0 baseline\\
SIGs Agent Sandbox \cite{sigAgent} &
Sandbox/Claim controller and optional WarmPool;
study reconciliation, allocation/refill and
declared idle-state behavior. &
P0\\
OpenSandbox BatchSandbox \cite{openSandbox} &
OpenSandbox API backed by Kubernetes Pods;
inspect the actual API/backend control path. &
P2 if distinct and feasible\\
OpenSandbox + Fast Sandbox \cite{fastSandbox} &
FastPath/Fastlet admission with multiple
Firecracker guests per Kubernetes-managed
worker; requires suitable host capabilities. &
P1 conditional\\
Agent Substrate \cite{substrate} &
Persistent actor identities, specialized
control plane, WorkerPool and checkpoint/
restore mechanisms; a task command endpoint
must be verified in the actor image. &
P1 conditional\\
\end{longtable}
}

The native cold baseline and any warm-Pod
allocator must be reported separately. Warm
capacity, pool refill, image caching,
resource reservations and preparation are
explicit controls. A native warm-pool
replica count does not establish a
platform's total concurrent-execution limit.

\subsection{State semantics and comparison strata}
Different lifecycle systems promise
different forms of state continuity:
\textbf{Tier 0} is an independently verified
command in an isolated environment;
\textbf{Tier 1} adds verified file/SQLite
continuity across the declared idle or
recreation path; and \textbf{Tier 2}
adds live process/memory restoration.

Compare all eligible systems on Tier 0,
and only systems that demonstrate the
relevant Tier 1 or Tier 2 contract on
those respective subtests.
Tier 2 is a separate optional
mechanism probe, not a prerequisite
for Scenario A/B or primary C1.
A PVC-backed Pod re-creation does
\emph{not} restore process memory.
An unsupported feature is reported
as N/A; installation failures
are reported separately from
task performance.

The extended five-platform capability
matrix and version-specific API
mapping are retained in the
technical appendix; documentation
alone does not establish that
a feature runs on Grid'5000.

\section{Benchmark Design and Methodology}
\label{sec:methodology}
\subsection{Application events and observable contract}
An external, deterministic driver maps application
events onto an adapter for each system.
The driver requires \texttt{acquire},
\texttt{execute} and \texttt{release},
and records passive state when available.
The optional \texttt{suspend}/\texttt{resume}
mechanism test (C2) is separate: the main
A and C1 workloads allow each platform
to follow its \emph{native} idle policy,
without a driver-forced state transition.

{\small
\begin{longtable}{@{}p{0.27\linewidth}p{0.61\linewidth}@{}}
\toprule
Event & Observable completion and distinction\\
\midrule
\endhead
\bottomrule
\endlastfoot
Acquire / first use &
Acquisition request until the \emph{first
independently verified fixed command}.
Record API response and Pod Ready,
when available, as intermediate milestones.\\
Execute &
Command submission until verified
exit status, output and required task files.\\
Idle / follow-up &
The idle gap ends on the new application
request. User-return latency is measured
from \emph{that request} until its first
verified command with correct retained state;
do not include the programmed idle gap.\\
Release &
Request until the logical identity is
inaccessible; separately observe physical
resource reclamation, reuse-ready capacity
and policy-retained data.\\
\end{longtable}
}

The workload is a pinned deterministic
application proxy with independent file,
test and/or database assertions. Live
LLM inference is kept off the timed path.
A successful API acknowledgement does
not establish task completion.

\subsection{Five core metric groups}\label{metric-definitions}
The Week 2 measurement dimensions are
retained in five groups, with no separate
mandatory platform-specific scorecard.

{\small
\begin{longtable}{@{}p{0.26\linewidth}p{0.48\linewidth}p{0.18\linewidth}@{}}
\toprule
Core metric & Primary quantities & Collection\\
\midrule
\endhead
\bottomrule
\endlastfoot
1. Lifecycle latency &
Request-to-first-verified-command, follow-up,
execution and logical cleanup latency;
per-run medians and raw small-cell samples;
p95 only with adequate sample size. &
Driver/verifier\\
2. Successful throughput &
Correctly completed tasks/trials or
return events per second versus
offered load; burst drain time. &
Driver events\\
3. Reliability and correctness &
Success, timeout, failure and incomplete
fractions; file/DB continuity, isolation
and expected final cleanup. &
Verifier/logs\\
4. Control-plane overhead &
Kubernetes API write requests and
Kubernetes/specialized management
CPU-seconds per verified unit;
publish raw totals. &
API/CPU metrics\\
5. Resource efficiency &
Runtime CPU-seconds, memory GiB-seconds
and active/idle cost; reserved
capacity reported separately. &
Cgroup/host\\
\end{longtable}
}

Optional \emph{bottleneck diagnostics}
include API latency histograms,
Scheduler pending Pods, controller
workqueues, etcd WAL fsync, pool
occupancy and snapshot I/O when
the deployed release actually exposes
them. A count of unfinished client
requests is \emph{not} internal queue
depth. Missing telemetry is N/A,
not zero.

\subsection{Comparability, normalization and statistics}
Pin the reserved cluster topology,
platform/Kubernetes versions,
task image and resources, driver
placement, cache state, network
conditions, preparation policy
and warm-pool budget. Differences
in isolation runtimes, storage or
prewarming that cannot be held
constant remain explicitly
reported confounding factors.

\textbf{Idle-policy configuration.} For every
platform and scenario, log the actual policy:
automatic pause/stop or reclaim enabled,
idle TTL/threshold, warm-pool size, retention
and cleanup grace period, and measured active
execution capacity. Record the observed native
transitions as well as the configured values.
Report default and explicitly tuned policies as
\emph{separate} configurations; do not pool their
results or attribute policy differences solely
to architecture. A/C1 never force a pause.

Compute API write counts (CREATE,
UPDATE, PATCH, DELETE) from
isolated-window counter changes,
and control-plane CPU consumption
from verified component counters.
API counters do \emph{not} label
each logical session: any
background-adjusted per-success
cost is an estimate. Keep all
failed attempts in the numerator,
publish raw totals and actual
success denominators, and mark
zero-success normalized cost
undefined. CPU-seconds and
GiB-seconds must not be added.

Plan five independent repetitions
for each feasible primary cell
(at least three if resources
prevent five). Reset logical
tasks and retain preparation
conditions and between-run
variability. \textbf{Small-sample rule:}
for fewer than 100 independently
verified completions in a
predeclared analysis cell, report
all client-observed latencies
(or the full sorted values),
median, range, completion counts
and failure/timeout rates.
Do not use a 5--10-request
P95/P99 as a stable tail result.
From 100 verified observations,
a descriptive p95 may be
reported with its sample count
and between-run variation; this
is a pragmatic minimum, not
a guarantee of statistical
precision. A p99 requires
substantially more samples
(e.g., at least 1,000
completed observations) and
is never mandatory. Preserve
per-run results even when a
pooled distribution is shown.
Success/timeout rates use
all offered requests; never
assign an invented completion
latency to timeouts.

Finally, distinguish
\emph{documented application
behavior}, measured trace data,
reproducible task proxies and
synthetic experimental parameters.
A shortened gap or synchronized
burst is a design choice, not
a production-arrival estimate.

\section{Use-case-grounded Scenarios}
\label{sec:scenarios}
These complementary workloads follow
\textbf{real application event patterns},
but their exact timing, batch size and
resource budgets are controlled experimental
parameters unless explicitly trace-derived.
\textbf{A and B are the primary study};
\textbf{C is optional Phase 2}.
Session, evaluation trial and persistent
user identity are different measurement units.

{\small
\begin{longtable}{@{}p{0.18\linewidth}p{0.36\linewidth}p{0.36\linewidth}@{}}
\toprule
Scenario & Application anchor & Distinct research question\\
\midrule
\endhead
\bottomrule
\endlastfoot
A (primary) &
Coding task followed by human review
and same-task continuation. &
How do idle and state reuse affect
the user's follow-up latency?\\
B (primary) &
Parallel short-lived evaluation trials
with automated verification. &
Where do provisioning, admission
and cleanup become bottlenecks?\\
C (optional) &
Independent persistent user identities
sharing constrained execution slots. &
How does native idle management
trade resource cost for return latency?\\
\end{longtable}
}

\subsection{A. Coding-agent continuation}
\label{scenario-a-coding-agent-continuation}

\textbf{Motivation.} Coding assistants can
edit a repository, run tools and tests,
then continue the same task after
developer feedback \cite{codexCloud,sweAgent}.
This tests the complete application session,
not only first Pod readiness.

\textbf{Workflow.} Acquire a sandbox for a
pinned repository and dependency image;
prepare a Python environment, apply
a deterministic first patch and verify tests;
wait through a review interval while
the platform applies its native idle policy;
apply a second patch to the \emph{same logical
task}, verify files/tests and release.
The scripted proxy has exactly
\textbf{two application turns and one
review gap}. A runtime may rebind its
physical container, but losing required
workspace state counts as failure.

\textbf{Parameters.} Use \textbf{10 primary
concurrent sessions}; 50 and 200 are
capacity-gated sensitivity levels.
A \textbf{120 s review gap} is synthetic,
not a replay of a production trace;
a separate 600 s native-policy probe
is optional. Provisional per-sandbox
resources are \textbf{2 vCPU,
4 GiB RAM, 10 GiB disk}. At 200
simultaneously reserved sandboxes,
the requested RAM alone would be
800 GiB before system overhead,
so this level must not be promised
on an unverified three-node testbed.

\textbf{Measurements.} First-verified-command
latency, follow-up request-to-verified-
command latency, file/test correctness,
idle CPU and memory cost, control-plane
work and logical cleanup. Observe actual
keep-alive, pause or recreation behavior
rather than assuming the gap triggers
eviction. Detailed command recipes and
the verifier are retained in the
technical appendix.

\subsection{B. Parallel agent evaluation}
\label{scenario-b-parallel-agent-evaluation}

\textbf{Motivation.} Evaluators and
agentic RL training systems issue
independent coding or terminal tasks
with multiple commands, verification
and environment cleanup. Sandman
and the official Harbor integration
provide concrete lifecycle anchors
\cite{sandman,terminalBench,harborIntegration}.

\textbf{Workflow.} A coordinator assigns
a batch of pinned trials. Each acquires
an isolated environment, runs 3--5
deterministic task steps or fixed oracle
commands, verifies the final result,
then releases resources. Trial state
does \emph{not} persist between
different evaluation tasks.

\textbf{Parameters.} Use a provisional
\textbf{1 vCPU/2 GiB} resource budget
and a \textbf{30--60 s} short-trial
target where feasible. Pilot at
\textbf{N=10/25}; start comparative
bursts at \textbf{N=50} only if
the pilot passes. Explore N=100,
200 and higher values subject
to actual cluster admission;
\textbf{N=1,000 is an optional
exploration ceiling}, not a
promised operating point. A
40/30/20/10 mix of I/O,
network/filesystem, CPU and
minimal tasks is an explicitly
synthetic heterogeneity option.

\textbf{Load modes.} The primary
synchronized-batch experiment
explores a \emph{degradation interval},
not one universal saturation knee.
For a low-sample exploratory sweep,
a provisional passing level requires
at least 95\% verified completions
before the declared deadline and
median first-command latency no
more than twice that platform's
low-load median. Repeated trials
then validate a separate
\emph{tail-latency degradation
criterion}: p95 at most twice
the low-load p95, only when
both baseline and evaluated
load have at least 100
successful observed first-command
latencies across declared
independent repetitions.
If samples are insufficient,
report a success/median-based
transition range and explicitly
mark the p95 boundary
\emph{undetermined}, rather than
claiming a validated tail knee.
Repeat the last passing and
first failing levels to
check verdict stability;
retain unresolved boundaries
and report the actual
hardware/admission limits.
The 95\% and $2\times$
thresholds are research
choices, not production SLOs.
A secondary sustained-churn
sweep at predeclared
$R_0/2$, $R_0$, and $2R_0$
holds task mix and
preparation fixed and
does not establish a precise
sustained saturation limit.

\textbf{Measurements.} Successful
verified trials/s versus offered
load, first-command and task
latency, failure/timeout rates,
API and specialized management
cost, worker utilization,
pool refill and cleanup.
The full load sweep and task
composition are archived in
the technical appendix.

\subsection{C. Persistent multi-user agents (optional Phase 2)}
\label{scenario-c-oversubscription-and-simultaneous-wake-phase-2}

\textbf{Motivation.} A covers
continued work within a coding
task; B covers short-lived trials.
Neither isolates \emph{many
persistent identities} sharing
a smaller verified execution
budget. Human-facing intermittent
agents and the documented Hermes
lifecycle motivate the pattern
\cite{dietz2025,datadog2026,hermesAaaS}.
They do not prove one universal
user-to-capacity ratio.

\textbf{Workflow.} Initialize
separate user identities with
private SQLite/files and verify
their initial commands. During
idle periods, the driver issues
no explicit suspend or resume.
User returns trigger a native
keep-alive, wake or recreation
path; independent assertions
must confirm the correct
\emph{user's} persisted state.

\textbf{Parameters.} The primary
design initializes \textbf{$N=20$
logical users}, a target
\textbf{$M=5$ usable execution
slots} ($R=N/M=4$),
\textbf{three requests per user}
separated by \textbf{two 300 s
idle gaps}, and about 1 MB
SQLite state plus two files
per identity. These are
\emph{synthetic design choices}.
The synthetic 300 s interval
is informed by observed
agent idle behavior
\cite{liu2026}; that study's
243 s median cross-turn
container resource-idle
interval is \emph{not}
a configured eviction
threshold.
A five-slot limit must be
verified on the actual platform;
it is not equivalent to
five warm-pool spare Pods.
$R=4$ alone does not prove
that requests will queue.

\textbf{Return modes.} The
\emph{primary} C1 test
staggers initial requests by
a nominal \textbf{10 s},
then schedules each user's
next request 300 s after its
previous verified command;
record actual inter-arrival
and idle times. A
\emph{secondary} one-second
synchronized stress probe
uses $B=5/10/20$ returning
users, with $B\leq N$.
Neither schedule is a
measured production arrival
distribution. Optional $N=50/100$,
$M=10/20$ and gap sensitivity
levels require verified
hardware capacity and
Tier 1 state support.
Explicit C2 suspend/restore
is an optional mechanism-only
test on eligible systems.

\textbf{Measurements.} Verified
per-user return latency,
success and state correctness,
idle runtime cost, resource
occupancy and observed native
lifecycle path. Admission
backlog is a diagnostic only
when the platform actually
reports queue events/metrics.
Staggered arrival results and
synchronized stress results
are always reported separately;
a small 5/10-user sample
must not drive a p99
performance conclusion.

\section{Experimental Setup and Feasibility}
\label{sec:experimental-setup}
\subsection{Grid'5000 testbed plan}
The initial candidate is the
Grid'5000 Rennes site, focusing
on \texttt{paradoxe} when an
appropriate reservation can be
obtained. The proposed topology
is one control-plane node and
two homogeneous Kubernetes
workers, but actual node
availability and the resource
allocation determine the setup.

For each run, record exact
node identifiers, CPU and
allocatable RAM, storage
and network topology,
kernel/cgroup mode, Kubernetes
and platform releases,
CNI/PVC configuration
and monitoring placement.
KVM/TUN and microVM snapshot
prerequisites are checked
only for relevant candidates;
the planned tasks include
CPU, filesystem, storage
and network activity, but
do not require a GPU.

\subsection{Pilot and shortlist}
The first objective is a
\emph{repeatable, observable
execution environment}, not
a peak-throughput claim.

\begin{enumerate}
\item Reserve suitable nodes,
deploy Kubernetes, and
verify core networking,
volume provisioning and
control-plane metrics.
\item Run a small preflight
task (write a file, perform
a computation, execute
a dependent command,
check output, then delete).
It is a \emph{smoke test},
not a fourth benchmark.
\item Complete the Tier 0
contract first on
\textbf{Native Pod} and
\textbf{SIGs Agent Sandbox};
validate Tier 1 only when
the scenario requires it.
\item Perform one small
fully logged A and B replay
before scaling; retain
failed and timed-out
requests and verify
physical cleanup.
\item Attempt Agent Substrate,
OpenSandbox Batch or Fast
only after runtime/storage
prerequisites and a
useful architecture
distinction are established.
\end{enumerate}

{\small
\begin{longtable}{@{}p{0.27\linewidth}p{0.60\linewidth}@{}}
\toprule
Pilot gate & Required evidence\\
\midrule
\endhead
\bottomrule
\endlastfoot
Functional &
Verified first and repeated
commands, output assertions
and cleanup of one session.\\
State &
File/DB continuity for
Tier 1; separate
process-memory evidence
if Tier 2 is claimed.\\
Capacity &
Actual Pod/worker admission
limit, image/warm-pool
budget and allocatable
host resources.\\
Telemetry &
Correlated driver timestamps,
Kubernetes API counters,
measured CPU/memory series
and identified missing
diagnostic signals.\\
Reproducibility &
Pinned versions, manifests,
task image, raw failures
and an archived run record.\\
\end{longtable}
}

A non-deployable candidate
is described as
\emph{infeasible on the
tested configuration},
not universally slower.
No Grid'5000 comparative
performance result is
established by this
planning document.

The complete OAR/Kadeploy
workflow, adapter endpoint
details, common-image
requirements and per-system
PoC record form are maintained
outside the proposal in
the technical appendix and
Grid'5000 deployment guide.

\section{Expected Analysis and Limitations}
\label{sec:analysis}
\subsection{Expected evidence and comparison}
First, compare \emph{verified
application outcomes} on
common feasible workload
cells: per-run latency,
successful throughput,
failure rates, measured
management work and
runtime resource efficiency.
Second, use optional
diagnostics to locate
symptoms in the Kubernetes
API/scheduler/controller,
specialized admission and
worker components,
guest runtime or storage.
Third, explain whether
warm preparation, native
idle policy, state semantics
or pool capacity change
the observed trade-off,
without assuming a
universal winner.

For Scenario B, report
both common-load comparisons
and the validated
\emph{degradation interval}.
For Scenario C, keep
staggered user experience,
secondary synchronized
stress and optional
explicit C2 mechanism
tests as distinct analyses.
Latency changes by
themselves do not prove
that a particular
component saturated.

\subsection{Threats to validity}
The application patterns
have documented precedents,
but the task proxies,
arrival schedules,
concurrency, resource
budgets and compressed
gaps are predominantly
\emph{synthetic experiment
choices}, not exact
production trace replays.
A small Grid'5000
cluster may expose
different bottlenecks
and absolute thresholds
from a larger fleet.
The same user-to-slot
ratio does not imply
equal queueing across
different deployments.

Cache state, snapshot
preparation, network
placement, CPU/guest
runtime, storage and
platform-native scheduling
cannot necessarily be
held identical; system
results should be
interpreted as
\emph{native-stack
comparisons with reported
co-varying factors}.
File continuity after
PVC-backed recreation
is distinct from
process-memory restoration.
Missing telemetry,
unsupported mechanisms
and installation failures
must be reported honestly.

Statistics are also
bounded by sample
counts, timeouts and
per-run variability:
a handful of returning
users cannot establish
reliable p99 tails.
Published vendor
figures and unverified
capacity assumptions
are not substitutes
for measurements on
the pinned releases.
The \href{https://github.com/HarvardMadSys/freeinference_agentic_trace}
{FreeInference Agentic Trace}
may inform time
distributions, but
does not contain
direct Kubernetes
lifecycle events
or prove individual
platform eviction
behavior.

\subsection{Planned deliverables and contribution boundary}
The proposed outputs
are (i) a reproducible
application-event driver
and platform adapters,
(ii) independent
task/state verifiers
and complete outcome
logs, (iii) five
core metric groups
at matched feasible
load points, and
(iv) qualified
cross-layer explanations
of architecture-specific
costs and degradation.
The practical priority
is the Native Pod
versus SIGs Agent
Sandbox comparison
on A/B; additional
architectures and C
require explicit
pilot feasibility.

The research contribution
is not simply which
product starts fastest,
but \emph{under which
workload conditions
and measured resource
budgets a lifecycle
strategy changes
user-visible completion,
reliability, management
work or idle cost}.
No result or
originality claim
is asserted before
the relevant experiments
and prior-work checks
are complete.

\clearpage
\phantomsection
\addcontentsline{toc}{section}{References}

\begin{thebibliography}{99}
\small

% References cited in the master proposal.
% Additional reading is indexed in docs/proposal-technical-appendix.md.

\bibitem{sandman}
Y. Yan, H. Wang, W. Chen, X. Wang, J. Liu, X. Li,
C. Qi, L. Fan, S. Li, Q. Wu, Y. Li, R. Wen, Q. Qu,
Z. Wang, and B. Dumoulin,
\emph{\href{https://arxiv.org/abs/2610.39711}
{Sandman: Building and Operating Sandboxes for Agentic
Reinforcement Learning on Shared Kubernetes}},
arXiv:2610.39711, 2026.

\bibitem{dietz2025}
M. Z. Pan et al.,
\emph{\href{https://arxiv.org/abs/2512.04123}
{Measuring Agents in Production}},
arXiv:2512.04123, 2026.
Accepted at ICML 2026 (Oral Presentation).

\bibitem{liu2026}
B. Liu, H. Qiu, I. Goiri, R. Fonseca,
R. Bianchini, and E. Choukse,
\emph{\href{https://arxiv.org/abs/2608.00101}
{Agentic Coding in the Wild: Characterizing
GitHub Copilot Traces at Production Scale}},
arXiv:2608.00101, 2026.

\bibitem{ariesPaper}
J. Park, L. Kondrashov, B. Zhou, H. Liu, Z. Liu,
C. Lu, R. Mancini, E. Choukse, H. Javaid,
G. Sviridov, T. Peng, C. Zhao, A. Avdeeva,
A. Gusev, M. Kogias, L. Mai, and D. Ustiugov,
\emph{\href{https://arxiv.org/abs/2607.29069}
{Rethinking AI Cloud Infrastructure for Agentic Serving
Systems with the Aries Experimentation Framework}},
arXiv:2607.29069, 2026.

\bibitem{deltaBox}
Y. Dong, J. He, S. Liu, Y. Hou, D. Du, Z. Xu,
S. Yu, B. Yang, Y. Xia, and H. Chen,
\emph{\href{https://arxiv.org/abs/2605.22781}
{DeltaBox: Scaling Stateful AI Agents with
Millisecond-Level Sandbox Checkpoint/Rollback}},
arXiv:2605.22781, 2026.

\bibitem{terminalBench}
M. A. Merrill et al.,
\emph{\href{https://arxiv.org/abs/2601.11868}
{Terminal-Bench: Benchmarking Agents on Hard,
Realistic Tasks in Command Line Interfaces}},
arXiv:2601.11868, 2026.

\bibitem{datadog2026}
C. Shamieh, A. Barksdale, J. Eastham, and O. Nzia,
\emph{\href{https://www.datadoghq.com/blog/serverless-ai-agents/}
{Default to Serverless for AI Agents, Then Let Production
Data Decide What Moves}},
Datadog Engineering Blog, October 29, 2026.

\bibitem{frigaBenchmark}
FrigaZzz,
\emph{\href{https://github.com/FrigaZzz/agent-sandbox-benchmark}
{AI Agent Sandbox Benchmark and Field Guide}},
Open-source benchmark repository, October 2026.

\bibitem{computeBenchmark}
ComputeSDK,
\emph{\href{https://github.com/computesdk/benchmarks}
{ComputeSDK Benchmarks}},
Open-source benchmark methodology and results repository,
accessed November 2026.

\bibitem{sigAgent}
Kubernetes SIG Apps,
\emph{\href{https://github.com/kubernetes-sigs/agent-sandbox}
{Agent Sandbox: Kubernetes-Native Sandboxes for Stateful
Workloads}},
Official source repository and documentation,
accessed November 2026.

\bibitem{openSandbox}
OpenSandbox Contributors,
\emph{\href{https://github.com/opensandbox-group/OpenSandbox}
{OpenSandbox}},
Official source repository and documentation,
accessed November 2026.

\bibitem{substrate}
Agent Substrate Contributors,
\emph{\href{https://github.com/agent-substrate/substrate}
{Agent Substrate: The Core System}},
Official source repository and architecture documentation,
accessed November 2026.

\bibitem{fastSandbox}
OpenSandbox Contributors,
\emph{\href{https://github.com/opensandbox-group/fast-sandbox}
{Fast Sandbox: High-Performance Sandbox Runtime on Kubernetes
with Firecracker microVMs}},
Official source repository and documentation,
accessed November 2026.

\bibitem{sigTuning}
Kubernetes SIG Apps,
\emph{\href{https://github.com/kubernetes-sigs/agent-sandbox/blob/main/docs/performance-tuning.md}
{Agent Sandbox Performance Tuning}},
Official engineering documentation,
accessed November 2026.

\bibitem{fastPerformance}
OpenSandbox Contributors,
\emph{\href{https://github.com/opensandbox-group/fast-sandbox}
{Fast Sandbox: Performance and Benchmark Documentation}},
ACK engineering measurements,
October 2026.

\bibitem{hermesAaaS}
Kubernetes SIG Apps,
\emph{\href{https://agent-sandbox.sigs.k8s.io/docs/use-cases/examples/hermes-agents-as-a-service/}
{Hermes Agents as a Service}},
Official use-case example,
accessed November 2026.

\bibitem{harborIntegration}
Kubernetes SIG Apps,
\emph{\href{https://agent-sandbox.sigs.k8s.io/docs/use-cases/examples/harbor-terminal-bench/}
{Terminal-Bench with Harbor}},
Official integration example,
accessed November 2026.

\bibitem{codexCloud}
OpenAI,
\emph{\href{https://help.openai.com/en/articles/20001545-using-codex-cloud}
{Using Codex Cloud}},
Official OpenAI Help Center documentation,
accessed November 2026.

\bibitem{sweAgent}
SWE-agent Contributors,
\emph{\href{https://github.com/SWE-agent/SWE-agent}
{SWE-agent: Agent-Computer Interfaces for Autonomous
Software Engineering}},
Open-source project repository,
accessed November 2026.

\bibitem{ariesRepo}
HyScale Lab,
\emph{\href{https://github.com/hyscale-lab/ARIES}
{ARIES: Agent Runtime and Infrastructure Experimentation System}},
Open-source research framework,
accessed November 2026.

\end{thebibliography}

\clearpage
\appendix
\section{Technical Appendix and Reproducibility Materials}
\label{sec:technical-appendix}

The proposal separates its research questions and methodology
from version-specific implementation notes. The companion
\href{https://github.com/Paperpaperdog/Agent_Sandbox_Proposal/blob/main/docs/proposal-technical-appendix.md}
{Technical Appendix (Markdown)} preserves task recipes,
adapter mappings, idle-policy and PoC checks, failure
taxonomy, archived reading and optional trace-calibration.
The \href{https://github.com/Paperpaperdog/Agent_Sandbox_Proposal/blob/924501eca5448bc4dca263364cbdb7f9e6718d04/natalia.tex}
{complete pre-consolidation proposal (commit 924501e)}
remains available for source-level audit.
Neither file constitutes measured benchmark results.

\end{document}
